% ===========================================================================
%  Annexe réservée au professeur — progression annuelle, sans dates
%  Incluse uniquement dans l'édition professeur (voir livre-t12s.tex).
% ===========================================================================
\chapter{Progression annuelle}
\markboth{Progression annuelle}{Progression annuelle}

Cette annexe ne porte aucune date. Le calendrier scolaire change chaque année,
les vacances se déplacent, une période est parfois amputée par les examens
blancs ou par les intempéries — mais la logique de la progression, elle, ne
bouge pas. Ce sont ces repères-là qui sont consignés ici : l'ordre des
chapitres, les volumes horaires du Programme d'Études, le découpage indicatif
en séances, et ce que l'élève doit savoir faire avant qu'on tourne la page.

Reportez vos propres dates dans la dernière colonne au début de l'année : le
tableau devient alors votre feuille de route.

\section*{Vue d'ensemble}

L'année officielle compte cinq périodes. Le livre, lui, est rangé par domaine.
Le tableau ci-dessous fait la conversion dans les deux sens.

\vspace{0.6em}
\begin{center}
\setlength{\tabcolsep}{6pt}
\renewcommand{\arraystretch}{1.5}
\begin{tabular}{@{}L{1.5cm} L{6.3cm} L{1.6cm} L{2.5cm}@{}}
\toprule
\textsf{\footnotesize\bfseries PÉRIODE} &
\textsf{\footnotesize\bfseries CHAPITRES} &
\textsf{\footnotesize\bfseries HEURES} &
\textsf{\footnotesize\bfseries VOS DATES} \\
\midrule
\textsf{\bfseries I}   & 1 Limites et continuité \newline 12 Nombres complexes \newline 2 Dérivabilité & 32 h & \\
\textsf{\bfseries II}  & 13 Barycentre \newline 14 Transformations (première moitié) \newline 3 Étude et représentation graphique \newline 4 Primitives et calcul intégral \newline 8 Suites numériques & 45 h & \\
\textsf{\bfseries III} & 5 Fonctions logarithmes \newline 6 Exponentielle et fonctions puissances \newline 7 Équations différentielles & 28 h & \\
\textsf{\bfseries IV}  & 14 Isométries (seconde moitié) \newline 15 Similitudes planes \newline 16 Géométrie dans l'espace \newline 9 Arithmétique & 39 h & \\
\textsf{\bfseries V}   & 10 Systèmes de numération \newline 11 Calcul matriciel \newline 17 Probabilités \newline 18 Variables aléatoires et lois \newline Révisions et épreuves blanches & 40 h & \\
\midrule
& \textsf{\bfseries Total} & \textsf{\bfseries 195 h} & \\
\bottomrule
\end{tabular}
\end{center}

\vspace{0.4em}
\noindent Le total correspond exactement aux volumes du Programme d'Études :
soixante heures pour l'analyse, quarante-cinq pour chacune des trois autres
composantes.

\section*{Découpage indicatif en séances}

Le découpage suppose des séances de deux heures. Adaptez-le : ce qui compte
est le rapport entre les chapitres, pas le nombre exact de séances.

\vspace{0.5em}
\begin{center}
\setlength{\tabcolsep}{6pt}
\renewcommand{\arraystretch}{1.35}
\begin{tabular}{@{}L{0.7cm} L{5.1cm} L{1.5cm} L{1.7cm} L{2.3cm}@{}}
\toprule
& \textsf{\footnotesize\bfseries CHAPITRE} &
\textsf{\footnotesize\bfseries HEURES} &
\textsf{\footnotesize\bfseries SÉANCES} &
\textsf{\footnotesize\bfseries ÉVALUATION} \\
\midrule
1  & Limites et continuité                     & 10 h & 5 & devoir court \\
2  & Dérivabilité                              & 10 h & 5 & devoir court \\
3  & Étude et représentation graphique         & 8 h  & 4 & — \\
4  & Primitives et calcul intégral             & 12 h & 6 & devoir surveillé \\
5  & Fonctions logarithmes                     & 10 h & 5 & — \\
6  & Exponentielle et fonctions puissances     & 10 h & 5 & devoir surveillé \\
7  & Équations différentielles                 & 8 h  & 4 & — \\
\midrule
8  & Suites numériques                         & 12 h & 6 & devoir surveillé \\
9  & Arithmétique                              & 14 h & 7 & devoir surveillé \\
10 & Systèmes de numération                    & 6 h  & 3 & — \\
11 & Calcul matriciel                          & 10 h & 5 & devoir court \\
\midrule
12 & Nombres complexes                         & 12 h & 6 & devoir surveillé \\
13 & Barycentre                                & 8 h  & 4 & — \\
14 & Transformations et isométries             & 10 h & 5 & devoir court \\
15 & Similitudes planes                        & 8 h  & 4 & — \\
16 & Géométrie dans l'espace                   & 10 h & 5 & devoir surveillé \\
\midrule
17 & Probabilités                              & 12 h & 6 & devoir court \\
18 & Variables aléatoires et lois              & 12 h & 6 & devoir surveillé \\
\bottomrule
\end{tabular}
\end{center}

\section*{Les repères de fin de période}

Ce que l'élève doit savoir faire avant de passer à la suite. Si un seul de ces
points n'est pas acquis, il vaut mieux reprendre une séance que d'avancer :
chacun de ces savoir-faire est réinvesti dès la période suivante.

\subsection*{À la fin de la première période}

Déterminer une limite, y compris celle d'une composée. Rédiger complètement
une application du théorème des valeurs intermédiaires — continuité, stricte
monotonie, encadrement, conclusion. Étudier la dérivabilité en un point et
reconnaître un point angulé. Passer d'une forme algébrique à une forme
trigonométrique d'un nombre complexe et résoudre une équation du second degré
dans $\C$.

\subsection*{À la fin de la deuxième période}

Construire un barycentre et déterminer une ligne de niveau. Mener une étude de
fonction jusqu'au tracé, asymptotes comprises. Reconnaître une forme
$u'u^{n}$ et calculer une intégrale par parties. Démontrer une propriété par
récurrence et établir la convergence d'une suite.

\subsection*{À la fin de la troisième période}

Résoudre équations, inéquations et systèmes faisant intervenir logarithme et
exponentielle, en posant les conditions d'existence avant tout calcul.
Reconnaître les limites usuelles des deux fonctions. Résoudre une équation
différentielle du premier et du second ordre, avec ou sans second membre.

\subsection*{À la fin de la quatrième période}

Décomposer une rotation en deux réflexions. Identifier une isométrie et une
similitude à partir de leur expression complexe. Écrire l'équation d'un plan
et discuter la position relative d'une droite et d'un plan. Résoudre une
équation diophantienne et manipuler les congruences.

\subsection*{À la fin de la cinquième période}

Effectuer des opérations matricielles et résoudre un système par les matrices.
Écrire un entier dans une base donnée. Construire un arbre de probabilité,
calculer une probabilité conditionnelle, donner la loi d'une variable
aléatoire et reconnaître une situation binomiale.

\section*{Les points de vigilance}

Trois passages coûtent régulièrement plus de temps que prévu, et il vaut mieux
l'avoir anticipé.

Le premier est l'arithmétique : quatorze heures paraissent confortables, mais
les équations diophantiennes et les congruences demandent une gymnastique de
raisonnement à laquelle les élèves ne sont pas préparés. Prévoyez une séance
de plus qu'annoncé, prise sur les révisions.

Le deuxième est la géométrie des transformations. La difficulté n'y est pas
technique : elle tient à la figure. Un élève qui trace mal ne trouve rien. Une
séance entière consacrée aux constructions à la règle et au compas, avant
toute théorie, est rentable.

Le troisième est la cinquième période, qui porte à la fois deux chapitres
neufs et l'ensemble des révisions. Elle est en pratique plus courte qu'elle
n'en a l'air. Si vous devez prendre de l'avance quelque part, prenez-la en
troisième période, où l'exponentielle avance vite une fois le logarithme
acquis.
